% 系统整体架构图（图 3.1）— 独立编译：xelatex system_architecture.tex
\documentclass[tikz,border=8pt]{standalone}
\usepackage{ctex}
\usetikzlibrary{arrows.meta,calc,positioning}

\begin{document}
\begin{tikzpicture}[
    font=\small,
    hdr/.style={
        rectangle,
        minimum width=4.35cm,
        minimum height=1.05cm,
        rounded corners=5pt,
        draw=black!55,
        very thick,
        align=center,
        font=\bfseries
    },
    card/.style={
        rectangle,
        minimum width=4.35cm,
        minimum height=3.45cm,
        rounded corners=6pt,
        draw=black!55,
        thick,
        align=left,
        inner sep=9pt,
        font=\scriptsize
    },
    flow/.style={
        -{Stealth[length=2.6mm]},
        very thick,
        draw=black!65
    },
    fb/.style={
        -{Stealth[length=2.6mm]},
        very thick,
        dashed,
        draw=red!68!black
    },
    lbl/.style={font=\scriptsize, inner sep=1pt, fill=white, fill opacity=0.92, text opacity=1}
]

% ---------- 三模块位置（横向流水线）----------
\coordinate (row) at (0,0);
\node[hdr, fill=orange!28, above=0pt of row, xshift=-5.35cm, anchor=south] (H1) {驱动函数接口\\静态分析与识别};
\node[card, fill=orange!9, below=0pt of H1.south, anchor=north] (C1) {
    \textbf{输入：}固件源码、CodeQL 数据库\\[0.35em]
    \textbullet\ MMIO 与驱动函数定位\\
    \textbullet\ 访问点、调用图、数据流\\
    \textbullet\ 统一数据模型 / 缓存查询接口
};

\node[hdr, fill=blue!22, draw=blue!65, above=0pt of row, anchor=south] (H2) {驱动函数分析、\\替换与编译};
\node[card, fill=blue!7, draw=blue!55, below=0pt of H2.south, anchor=north] (C2) {
    \textbf{核心智能层（LangGraph）}\\[0.35em]
    \textbullet\ Analyzer Agent：分类与替换生成\\
    \textbullet\ VerifyReplacement 定义校验\\
    \textbullet\ 写入源码、工程编译（Build）\\
    \textbullet\ BuilderFixer：编译失败修复闭环
};

\node[hdr, fill=green!24, above=0pt of row, xshift=5.35cm, anchor=south] (H3) {动态测试\\反馈闭环};
\node[card, fill=green!8, below=0pt of H3.south, anchor=north] (C3) {
    \textbf{基于仿真器（如 Unicorn）}\\[0.35em]
    \textbullet\ 运行替换后固件\\
    \textbullet\ 日志、崩溃、覆盖率采集\\
    \textbullet\ 结构化错误 / 失败样例\\
    \textbullet\ 回流触发再分析与再生成
};

% ---------- 正向数据流 ----------
\draw[flow] (C1.east) -- node[lbl, above=2pt] {驱动函数列表与上下文} (C2.west);
\draw[flow] (C2.east) -- node[lbl, above=2pt] {替换代码 / 可构建镜像} (C3.west);

% ---------- 反馈回路（经底部绕行）----------
\path let \p1=(C3.south), \p2=(C2.south) in coordinate (FbR) at (\x1, \y1-0.95)
    coordinate (FbM) at (\x2, \y1-0.95);
\draw[fb] (C3.south) -- (FbR) -- node[lbl, below=2pt] {失败样例与诊断信息} (FbM) -- (C2.south);

% ---------- 底部总览 ----------
\node[below=1.35cm of C2.south, font=\footnotesize\bfseries, text=gray!75!black, align=center] {
    静态识别 $\rightarrow$ 智能替换与构建 $\rightarrow$ 动态验证 \quad|\quad 反馈闭环迭代优化
};

\end{tikzpicture}
\end{document}
